\documentclass[12pt,a4paper,oneside]{extreport}

% Pacotes essenciais
\usepackage{graphicx}

% Fontes
\usepackage{fontspec}
\setmainfont{TeX Gyre Pagella}

% Idioma
\usepackage[brazil]{babel}

% Margens
\usepackage[
    left=3cm,
    right=2.5cm,
    top=2.5cm,
    bottom=3cm
]{geometry}

% Espaçamento
\usepackage{setspace}
\onehalfspacing

% Microtipografia
\usepackage{microtype}

% Cabeçalhos e Rodapés
\usepackage{fancyhdr}
\setlength{\headheight}{16pt}
\pagestyle{fancy}
\fancyhf{} % Limpa
\fancyhead[R]{{\leftmark}}
\fancyhead[L]{\textbf{Apostila de Programação C}}
\renewcommand{\headrulewidth}{0.4pt}
\renewcommand{\footrulewidth}{0pt}

% Primeira página de cada capítulo
\fancypagestyle{plain}{
    \fancyhf{}
    \fancyfoot[C]{\thepage}
    \renewcommand{\headrulewidth}{0pt}
}

% Hyperref
\usepackage[hidelinks]{hyperref}

% Pacote de Códigos
\usepackage{listings}
\lstset{
    language=C,
    basicstyle=\ttfamily\small,
    keywordstyle=\color{blue},
    commentstyle=\color{green!60!black},
    stringstyle=\color{red},
    numbers=left,
    numberstyle=\tiny,
    frame=single,
    breaklines=true
}

% Cores
\usepackage{xcolor}
\definecolor{titulo}{RGB}{20,55,120}

% Caixas (tcolorbox)
\usepackage[most]{tcolorbox}

% Definições de Caixas Pedagógicas
\newtcolorbox{observacao}[1][]{
    colback=blue!5!white,
    colframe=blue!75!black,
    title=Observação,
    fonttitle=\bfseries,
    #1
}

\newtcolorbox{importante}[1][]{
    colback=red!5!white,
    colframe=red!75!black,
    title=Importante,
    fonttitle=\bfseries,
    #1
}

\newtcolorbox{curiosidade}[1][]{
    colback=green!5!white,
    colframe=green!50!black,
    title=Curiosidade,
    fonttitle=\bfseries,
    #1
}

\newtcolorbox{errocomum}[1][]{
    colback=orange!5!white,
    colframe=orange!80!black,
    title=Erro Comum,
    fonttitle=\bfseries,
    #1
}

\newtcolorbox{objetivos}[1][]{
    colback=gray!5!white,
    colframe=gray!50!black,
    title=Objetivos de Aprendizagem,
    fonttitle=\bfseries,
    #1
}

\newtcolorbox{resumo}[1][]{
    colback=yellow!5!white,
    colframe=yellow!60!black,
    title=Resumo,
    fonttitle=\bfseries,
    #1
}

\newtcolorbox{exercicio}[1][]{
    colback=purple!5!white,
    colframe=purple!50!black,
    title=Exercício,
    fonttitle=\bfseries,
    #1
}

\newtcolorbox{desafio}[1][]{
    colback=black!5!white,
    colframe=black!80!white,
    title=Desafio,
    fonttitle=\bfseries,
    #1
}

\newtcolorbox{terminal}[1][]{
    colback=black!90!white,
    colframe=black,
    coltext=white,
    fontupper=\ttfamily,
    title=Terminal,
    fonttitle=\bfseries,
    #1
}

% Títulos
\usepackage{titlesec}
\titleformat{\chapter}
{\Huge\bfseries\color{titulo}}
{\thechapter}
{1em}
{}

\setlength{\parskip}{0.4em}


\title{\huge Apostila de Preparação para Programação \textbf{C}}
\author{Rafael Coelho}
%\date{July 2026}

\begin{document}
\maketitle

\chapter*{Prefácio}
\addcontentsline{toc}{chapter}{Prefácio}

A linguagem C ocupa um papel fundamental na formação de profissionais da Computação. Embora linguagens modernas ofereçam mecanismos que abstraem diversos detalhes da implementação, compreender como um programa é executado, como os dados são organizados na memória e como os recursos do computador são utilizados continua sendo uma habilidade indispensável para quem deseja construir uma base sólida em programação.

Esta apostila foi elaborada como material de apoio às monitorias da disciplina de Linguagem C. Seu propósito não é substituir as aulas ou os livros clássicos da área, mas complementar o processo de aprendizagem por meio de explicações objetivas, exemplos comentados, diagramas de memória, exercícios progressivos e observações sobre os erros mais frequentes encontrados pelos estudantes.

Ao longo dos capítulos, a linguagem será apresentada não apenas como uma ferramenta para escrever código, mas como um meio para compreender os fundamentos da programação e o funcionamento interno dos computadores. Sempre que possível, os conceitos serão relacionados a situações práticas e acompanhados de representações visuais que auxiliem na construção de um modelo mental consistente sobre a execução dos programas.

Espera-se que este material sirva como um roteiro de estudos para alunos que desejam consolidar os conhecimentos adquiridos em sala de aula, preparar-se para avaliações e, principalmente, desenvolver uma compreensão mais profunda dos princípios que fundamentam a programação em baixo nível.

\vspace{1cm}

\begin{flushright}
\textit{Bom estudo!}
\end{flushright}

\chapter*{Como utilizar esta apostila}
\addcontentsline{toc}{chapter}{Como utilizar esta apostila}

\tableofcontents
\newpage

\part{Fundamentos de C}

\input{capitulos/01-introducao.tex}
\input{capitulos/02-variaveis-tipos-operadores.tex}

% TODO: Include remaining chapters as they are developed
%\input{capitulos/03-controle-fluxo.tex}
%\input{capitulos/04-funcoes.tex}

\part{Organização dos Dados}
%\input{capitulos/05-vetores.tex}
%\input{capitulos/06-strings.tex}

\part{Gerenciamento de Memória}
%\input{capitulos/07-structs-enums.tex}
%\input{capitulos/08-ponteiros.tex}
%\input{capitulos/09-alocacao-dinamica.tex}

\part{Programação Modular e Avançada}
%\input{capitulos/10-modularizacao.tex}
%\input{capitulos/11-estruturas-dados.tex}
%\input{capitulos/12-recursao.tex}
%\input{capitulos/13-bit-a-bit.tex}
%\input{capitulos/14-depuracao.tex}

\appendix
\part{Apêndices}
\chapter{Perguntas Frequentes}
\chapter{Erros Clássicos em C}

\end{document}
